\documentclass[11pt,a4paper]{article}

% ---------- Codifica e lingua ----------
\usepackage[T1]{fontenc}
\usepackage[utf8]{inputenc}
\usepackage[main=italian,provide=*]{babel}

% ---------- Matematica ----------
\usepackage{amsmath,amssymb}
\usepackage{mathtools}

% ---------- Impaginazione ----------
\usepackage[a4paper,margin=2.7cm]{geometry}
\usepackage{caption}
\usepackage{booktabs}
\usepackage{enumitem}
\usepackage{placeins}
\usepackage[colorlinks=true,linkcolor=black,citecolor=black,urlcolor=blue!55!black]{hyperref}

\newcommand{\ket}[1]{\lvert #1 \rangle}
\newcommand{\bra}[1]{\langle #1 \rvert}

\title{Rumore quantistico sul dimero\\
\large Panoramica della pipeline (vista a 2000 piedi)}
\author{Samuele \\ Relatore: Prof.\ Alessandro Chiesa}
\date{}

\begin{document}
\maketitle

\begin{abstract}
\noindent Questo documento è una mappa, non una derivazione: serve a
vedere in un colpo d'occhio come è strutturata la pipeline di rumore sul
dimero prima di entrare nel dettaglio di ciascuno stadio (ciascuno ha un
proprio documento dedicato, elencato in Sez.~\ref{sec:stadi}). Scope,
canali di rumore e decisioni operative sono quelli confermati dal
relatore (risposta del 2 settembre 2026, \texttt{domande\_relatore.md},
sez.~12) e seguono la terminologia delle sue slide
(\texttt{7Physical\_implementation.pdf}).
\end{abstract}

\tableofcontents

\section{Da dove partiamo}

La parte a sistema chiuso è completa sul dimero: Hamiltoniana esatta,
VQE con termine DM, evoluzione temporale $e^{-iHt}$ via
Suzuki--Trotter, e correlazioni dinamiche $C_{ij}^{\alpha\beta}(t)$
tramite un circuito di tipo Hadamard test. Tutto lavora su uno stato
\emph{puro} $\ket{\psi}$ che evolve unitariamente: non c'è interazione
con l'ambiente, non c'è rumore.

Questa pipeline ripercorre esattamente la stessa catena di stadi ---
preparazione dello stato fondamentale, evoluzione temporale, misura dei
correlatori --- ma la accompagna con due sorgenti di rumore realistiche,
quelle di un dispositivo quantistico vero. Non si introduce fisica nuova
nell'Hamiltoniana: si aggiunge un modello di come un computer quantistico
reale esegue (in modo imperfetto) gli stessi circuiti già scritti.

\section{Cosa cambia: da \texorpdfstring{$\ket{\psi}$}{|psi>} a \texorpdfstring{$\rho$}{rho}}

Il rumore accoppia il sistema a un ambiente che non osserviamo. Uno
stato puro non basta più a descrivere il sistema una volta che
tracciamo via i gradi di libertà dell'ambiente: serve l'operatore
densità $\rho$, e l'evoluzione non è più unitaria pura ma un'operazione
quantistica più generale. Il documento sul modello di rumore
(\texttt{risultati\_modello\_\allowbreak rumore\_dimero.tex}) richiama questo
formalismo dal punto dove serve; qui basta sapere che da questo momento
lavora con $\rho$ al posto di $\ket{\psi}$, e che in Qiskit questo
corrisponde a usare
\texttt{AerSimulator(method=\allowbreak"density\_matrix")} invece del
simulatore statevector.

\section{Gli stadi della pipeline}
\label{sec:stadi}

\begin{table}[htbp]
\centering
\begin{tabular}{@{}cll@{}}
\toprule
\# & Stadio & Cosa cambia rispetto al sistema chiuso \\
\midrule
1 & Modello di rumore & nuovo: nessun analogo nel sistema chiuso \\
2 & VQE con termine DM, rumoroso & stesso ansatz, simulatore a $\rho$ \\
3 & Quantum simulation (Trotter), rumorosa & stesso circuito, compilato per passo \\
4 & Correlazioni dinamiche, rumorose & stesso Hadamard test, readout incluso \\
5 & Scan sui parametri di rumore & nuovo: nessun analogo nel sistema chiuso \\
\bottomrule
\end{tabular}
\caption{I cinque stadi della pipeline. Gli stadi 2--4 sono mirror dei
corrispondenti stadi a sistema chiuso; gli stadi 1 e 5 sono propri di
questa pipeline.}
\end{table}
\FloatBarrier

\subsection{Modello di rumore}

Il relatore ha confermato due soli canali (slide 17 del corso): errore
di gate (canale depolarizzante, separato a 1 e 2 qubit) ed errore di
lettura (readout). $T_1$/$T_2$ sono esclusi esplicitamente. Il compito
di questo stadio è tradurre dati di calibrazione reali --- pubblicati
per un chip esistente --- in un oggetto \texttt{NoiseModel} di Qiskit,
passando dal formalismo a operatori di Kraus richiamato a lezione.
Documento: \texttt{risultati\_modello\_rumore\_dimero.tex} (con
derivazione completa in \texttt{teoria\_modello\_rumore.tex}).

\subsection{VQE con termine DM, sotto rumore}

Si ripete la stima del ground state con lo stesso ansatz PMA-2q$\cdot$3
già validato a sistema chiuso (fedeltà 1 nel caso ideale), ma il
circuito viene ora eseguito su un simulatore a operatore densità con il
\texttt{NoiseModel} del modello di rumore agganciato. L'osservabile
d'interesse è quanto l'energia stimata e la fedeltà rispetto al ground
state esatto si allontanano dal caso ideale. Documento:
\texttt{risultati\_vqe\_dm\_rumoroso\_dimero.tex}.

\subsection{Quantum simulation (Trotter), sotto rumore}
\label{subsec:trotter}

Stesso circuito di Suzuki--Trotter del sistema chiuso. Qui c'è
un'insidia tecnica verificata esplicitamente: se si transpila
\emph{l'intero} circuito a un livello di ottimizzazione alto, Qiskit
riconosce che gli $N$ passi identici formano nel complesso un unico
unitario a 2 qubit e li collassa in un circuito a costo
\emph{costante}, indipendente da $N$. Questo distruggerebbe
silenziosamente proprio il compromesso che vogliamo studiare (più passi
$\Rightarrow$ meno errore di Trotter ma più rumore accumulato). La
soluzione è transpilare una volta il singolo passo al suo costo minimo
(3 CNOT, già dimostrato in
\texttt{circuito\_compatto\_teoria\_e\_limiti.tex}) e poi ripetere il
blocco già compilato $N$ volte, senza ritranspilare il circuito
completo. Documento: \texttt{risultati\_trotter\_rumoroso\_dimero.tex}
--- include anche un cross-check indipendente (circuito Qiskit contro
pura potenza di matrice in \texttt{numpy}) e il confronto diretto fra
$F(N)$ a rumore nullo e a rumore acceso, per rendere esplicito perché
esista un compromesso.

\subsection{Correlazioni dinamiche, sotto rumore}

Stesso circuito Hadamard test del sistema chiuso (ancilla + 2 qubit di
registro), ora con rumore su tutti i suoi ingredienti: preparazione
dello stato, blocco $U(t)$ ripetuto (stadio precedente), e infine la
misura sull'ancilla, che include anche l'errore di lettura del modello
di rumore. Documento:
\texttt{risultati\_correlatori\_rumorosi\_dimero.tex}.

\subsection{Scan sui parametri di rumore}

Il relatore chiede esplicitamente di non fermarsi a un solo punto: si
ripete il conto dello stadio precedente (o già di quello di Trotter) al
variare di $(\varepsilon_{1q},\varepsilon_{2q},p_\text{readout})$, per
vedere come si sposta il numero di passi ottimale $N^*$ che minimizza
l'errore totale --- risultato di un compromesso fra errore di Trotter
(diminuisce con $N$) e rumore accumulato (cresce con $N$). Documento:
\texttt{risultati\_scan\_parametri\_rumore\_dimero.tex} --- include anche
un affinamento sulla legge di scala di $N^*$ verso la saturazione.

\section{Cosa resta fuori}

Per completezza, dichiarato esplicitamente per non lasciarlo implicito:
rilassamento e defasamento ($T_1$, $T_2$) sono esclusi per decisione
del relatore, anche se compaiono nelle sue slide come canali standard;
il readout asimmetrico è opzionale, trattato come estensione separata
(\texttt{teoria\_readout\_asimmetrico.tex},
\texttt{risultati\_readout\_asimmetrico.tex}) --- lo stesso vale per la
domanda se rioptimizzare il VQE sotto rumore cambi i risultati a valle
(\texttt{risultati\_vqe\_noise\_aware\_definitivo.tex}); le reti neurali
(NQS) restano fuori scope come nel sistema chiuso.

\section{Il filo conduttore verificato in ogni stadio}

Una proprietà usata, esplicitamente o implicitamente, in ogni stadio a
valle del modello di rumore: il readout simmetrico è sempre un fattore
\textbf{moltiplicativo} sull'osservabile misurato, mai additivo --- da
cui discende, per esempio, che $N^*$ non dipende da $p_\text{readout}$
(stadio 5). Questa proprietà smette di valere non appena il readout
diventa asimmetrico (\texttt{teoria\_readout\_\allowbreak asimmetrico.tex}): un
promemoria che ogni argomento di invarianza in questa pipeline poggia su
un'ipotesi verificata, non su una coincidenza.

\end{document}
